\section{Métodos Modernos de la Construcción}

Los Métodos Modernos de Construcción (MMC) son un conjunto de métodos mediante los cuales se aplica la Construcción Industrializada (CI) en un proyecto, con el objetivo de mejorar la productividad y la sustentabilidad a lo largo del ciclo de vida de la obra \cite{GuiaMMC}. La guía de CTEC y la Cámara Chilena de la Construcción organiza estas soluciones en siete categorías, desde módulos volumétricos fabricados fuera de obra (MMC 1) hasta tecnologías aplicadas directamente en sitio (MMC 7). A partir de lo observado en la visita al Parque CTEC, se identifica una oportunidad concreta de aplicación dentro de este proyecto.

\subsection{Identificación de la oportunidad y aplicación}

La oportunidad más clara es el campamento de los trabajadores, que se resuelve con un campamento modular prefabricado, correspondiente a la categoría MMC 1: Módulos Estructurales 3D (elementos volumétricos espaciales que forman parte del sistema estructural y se fabrican en condiciones controladas) \cite{GuiaMMC}.

A diferencia de las estructuras de tratamiento, que por su geometría circular, carácter monolítico y exigencia de impermeabilidad deben ejecutarse in situ, el campamento es un conjunto de recintos de geometría modular, repetitiva y de baja complejidad estructural, condiciones idóneas para su industrialización. El emplazamiento refuerza esta elección: el sector de La Negra es una zona desértica y remota donde la logística de materiales y mano de obra es costosa y las condiciones ambientales dificultan las faenas en terreno.

Los módulos se fabricarían en planta en ambiente controlado, en la modalidad de elemento terminado con estructura, aislación, revestimientos, instalaciones, puertas y ventanas ya integrados. Luego se transportan y montan con grúa sobre la losa de fundación de hormigón armado descrita en la Sección~\ref{sec:campamento}, que actúa como la única interfaz ejecutada in situ.
\subsection{Justificación}

La elección de la prefabricación modular se justifica por los siguientes beneficios, alineados con los indicadores reportados para proyectos industrializados en Chile \cite{GuiaMMC}:

\begin{itemize}
    \item Reducción de plazos y productividad: la fabricación en paralelo con las faenas de obra gruesa acorta significativamente los plazos. Estudios nacionales reportan disminuciones de plazo de 40 a 75\% y aumentos de productividad de 48 a 127\% respecto de la construcción tradicional.
    \item Menor mano de obra en terreno: la fabricación fuera de sitio reduce las horas-hombre en obra en 70 a 90\%, aspecto especialmente valioso en una zona remota donde alojar y movilizar personal es costoso.
    \item Sustentabilidad: la producción en ambiente controlado reduce los residuos de construcción en 65 a 70\%, y la solución modular permite el desmontaje y reutilización al término del proyecto, favoreciendo un esquema de economía circular.
    \item Seguridad laboral e independencia climática: la industrialización reduce la exposición a faenas de riesgo en terreno, y al fabricarse en planta los módulos no se ven afectados por la alta radiación, los vientos ni las variaciones térmicas del desierto.
\end{itemize}

En consecuencia, el campamento de trabajadores es la partida del proyecto donde la aplicación de un MMC genera el mayor valor agregado en productividad, plazos y sustentabilidad, sin comprometer las exigencias que obligan a mantener las estructuras de tratamiento en una solución tradicional in situ.
